\section{Komplekse tal}
\subsubsection{Definition}
\begin{align*}
    \sqrt{-1} &= i \\
    \sqrt{-n} &= i \cdot \sqrt{n} \quad  n,x \in \mathbb{R}^+ \\
    \\
    z &= a + bi \quad a,b \in \mathbb{R} \\
    \Re(z) &= a \\
    \Im(z) &= b
\end{align*}

\subsection{Aritmetik}
\subsubsection{Addition og Subtraktion}
\begin{align*}
    z_1 + z_2 &= (a + bi) + (c + di) = (a+c) + (b+d)i \\
    z_1 - z_2 &= (a + bi) - (c + di) = (a-c) + (b-d)i \\
\end{align*}

\subsubsection{Multiplikation}
\[
    z_1 \cdot z_2 = (a + bi)(c + di) = (a \cdot c - b \cdot d) + (a \cdot d + b \cdot c)i
\]

\paragraph{Særtilfælde: Multiplikation med det komplekst konjugerede}
\[
    (a+bi) \cdot (a-bi) = a^2 + b^2 \kildetag{Lemma 4.2.1} \\
\]


\subsubsection{Division}
\[
    \frac{z_1}{z_2} = \frac{a + bi}{c + di} = \frac{(a + bi)(c - di)}{(c + di)(c - di)} = \frac{(a \cdot c + b \cdot d) + (b \cdot c - a \cdot d)i}{c^2 + d^2}
\]

\paragraph{Særtilfælde: 1 over et komplekst tal}
\[
    \frac{1}{c + di} = \frac{c - di}{(c + di)(c - di)} = \frac{c}{c^2 + d^2} - \frac{d}{c^2 + d^2}i \kildetag{Ligning 4.3} \\
\]


\subsection{Kompleks Konjugering}

Det komplekst konjugerede af $z=a+bi$ er $\overline{z}=a-bi$.
\begin{align*}
    (z) \cdot (\overline{z}) &= {\Re(z)}^2 + {\Im(z)}^2 \kildetag{Lemma 4.2.1} \\
    \overline{\overline{z}} &= z \kildetag{Lemma 5.3.1} \\
    \overline{z_1 + z_2} &= \overline{z_1} + \overline{z_2} \kildetag{Lemma 5.3.1} \\
    \overline{z_1 \cdot z_2} &= \overline{z_1} \cdot \overline{z_2} \kildetag{Lemma 5.3.1} \\
    \overline{\left( \frac{1}{z} \right)} &= \frac{1}{\overline{(z)}} \kildetag{Lemma 5.3.1} \\
    \overline{z^n} &= (\overline{z})^n \quad n \in \mathbb{Z} \kildetag{Lemma 5.3.1} \\
    \overline{e^z} &= e^{\overline{z}} \kildetag{Lemma 5.3.2} \\
    \overline{e^{ia}} &= e^{-ia} \quad a \in \mathbb{R} \kildetag{Lemma 5.3.2} \\
    \overline{z} &= |z| \cdot e^{-i \cdot \arg(z)} \kildetag{Lemma 5.3.2} \\
\end{align*}

\bemærkBig{Bemærk \(|e| \cdot e^{-i \cdot \arg(z)}\) er den polære form for det komplekst konjugerede af \(z\).}

\subsection{Modulus og argument}

\begin{center}
\begin{minipage}{0.4\textwidth}
\begin{align*}
    r &= |z| = \sqrt{a^2 + b^2} \\
    \theta &= \arg(z) = \tan^{-1}\left(\frac{b}{a}\right)
\end{align*}
\end{minipage}
\hfill
\begin{minipage}{0.5\textwidth}
\begin{tikzpicture}[scale=2, >=Stealth]
    % Axes
    \draw[->] (-0.2,0) -- (3,0);
    \draw[->] (0,-0.2) -- (0,2.2);
    
    % Point z
    \coordinate (O) at (0,0) node[above left] {$0$};
    \coordinate (z) at (2,1.5);
    \coordinate (xaxis) at (3,0);
    \filldraw (z) circle (1.5pt) node[above right] {$z$};
    
    % Dotted line from origin to z
    \draw[dotted] (O) -- (z) node[midway, above left] {$|z|$};
    
    % Projection lines
    \draw[dashed] (z) -- (z |- O) node[midway, right] {$\Im(z)$};
    \draw[dashed] (z) -- (z -| O) node[midway, above] {$\Re(z)$};
    
    % Angle
    \pic[draw, ->, angle radius=1cm, angle eccentricity=1.2] 
        {angle = xaxis--O--z};
    \node at (1.1,0.25) {arg$(z)$};
\end{tikzpicture}
\end{minipage}
\end{center}

For ethvert argument $\theta$ vil $\theta + 2k\pi$ for $k \in \mathbb{Z}$ være det samme argument. 

Argumentet i intervallet $]-\pi, \pi]$ kaldes hovedargumentet og skrives $\Arg(z)$.

\subsubsection{Trigonometri}
\begin{align*}
    \Re(z) &= |z| \cdot \cos(\arg(z)) \kildetag{Ligning 4.4} \\
    \Im(z) &= |z| \cdot \sin(\arg(z)) \kildetag{Ligning 4.4} \\
    z &= |z| \cdot \left(\cos(\arg(z)) + \sin(\arg(z)) \cdot i \right) \kildetag{Ligning 4.5} \\
    |z| &= \sqrt{\Re(z)^2 + \Im(z)^2} \kildetag{Ligning 4.6} \\
\end{align*}

\subsubsection{Find argument of modulus fra rektangulær form}
\[
    |z| = \sqrt{a^2 + b^2} \kildetag{Sætning 4.3.1}
\]

\begin{center}

\begin{minipage}{0.55\textwidth}
\[
    \Arg(z) = \begin{cases}
        \arctan\left(\frac{b}{a}\right) & a > 0 \\
        \frac{\pi}{2} & a = 0, b > 0 \\
        \arctan\left(\frac{b}{a}\right) + \pi & a < 0, b \geq 0 \\
        -\frac{\pi}{2} & a = 0, b < 0 \\
        \arctan\left(\frac{b}{a}\right) - \pi & a < 0, b < 0 \\
    \end{cases} 
\]
\end{minipage}
\hfill
\begin{minipage}{0.4\textwidth}
\begin{tikzpicture}[scale=1.5, >=stealth]
    % Axes
    \draw[->] (-2,0) -- (2,0) node[below] {$a$};
    \draw[->] (0,-2) -- (0,2) node[left] {$b$};
    
    % Quadrant labels and cases
    \node at (1,0.7) {$\arctan\left(\frac{b}{a}\right)$};
    \node at (-1,0.7) {$\arctan\left(\frac{b}{a}\right) + \pi$};
    \node at (-1,-0.7) {$\arctan\left(\frac{b}{a}\right) - \pi$};
    \node at (1,-0.7) {$\arctan\left(\frac{b}{a}\right)$};
    
    % Special cases on axes with white backgrounds
    \node[above, fill=white, inner sep=2pt] at (0,1) {$\frac{\pi}{2}$};
    \node[below, fill=white, inner sep=2pt] at (-0.1,-1) {$-\frac{\pi}{2}$};
\end{tikzpicture}
\end{minipage}
\end{center}

\kilde{Sætning 4.3.1}

\subsubsection{Aritmetik}
Lad $z_1, z_2 \in \mathbb{C} \backslash \{ 0 \}$, så gælder:
\begin{align*}
    |z_1 \cdot z_2| &= |z_1| \cdot |z_2| \\
    \arg(z_1 \cdot z_2) &= \arg(z_1) + \arg(z_2) \\
    |z_1 / z_2| &= |z_1| / |z_2| \\
    \arg(z_1 / z_2) &= \arg(z_1) - \arg(z_2) \\
    |z^n| &= |z|^n \quad n \in \mathbb{Z} \\
    \arg(z^n) &= n \cdot \arg(z) \quad n \in \mathbb{Z}
\end{align*}
\kilde{Sætning 4.6.2}


\subsection{Den komplekse eksponentialfunktion}
\[
    e^z = e^{a + bi} = e^a \cdot e^{bi} = e^a \cdot \left(\cos(b) + \sin(b) \cdot i \right) \kildetag{Definition 4.4.1}
\]

Den komplekse eksponentialfunktion betegnes ofte som \(\exp(t)= e^t\)


\subsection{Polære koordinater}
\bemærk{Se "Polær form" længere nede}

Talparret $(|z|, \Arg(z))$ kaldes polære koordinater for et komplekst tal $z$.

Hvis et komplekst tal er skrevet \(z = r \cdot \left(\cos(\alpha) + \sin(\alpha) \cdot i \right) \) hvor $r$ er et positivt reelt tal, så er $|z|=r$ og $\arg(z)=\alpha$.


\subsection{Yderligere}
\[
    (\cos(\alpha_1) + \sin(\alpha_1) i) \cdot (\cos(\alpha_2) + \sin(\alpha_2) i) = \cos(\alpha_1 + \alpha_2) + \sin(\alpha_1 + \alpha_2) i \\
\]
\kilde{Lemma 4.4.1}

\begin{align*}
    e^z &\neq 0 \\
    \frac{1}{e^z} &= e^{-z} \\ 
    e^{z_1} \cdot e^{z_2} &= e^{z_1 + z_2} \\
    e^{z_1}/e^{z_2} &= e^{z_1 - z_2} \\
    (e^{z})^n &= e^{n \cdot z} \quad n \in \mathbb{Z}
\end{align*}
\kilde{Sætning 4.4.2}

\subsection{Eulers formel}
\[
    e^{it} = \cos(t) + i \sin(t) \kildetag{Euler's formel}
\]
Medfører:
\begin{align*}
    e^{-it} &= \cos(t) - i \sin(t) \kildetag{Ligning 4.8} \\
    \cos(t) &= \frac{e^{it} + e^{-it}}{2} \kildetag{Ligning 4.9} \\
    \sin(t) &= \frac{e^{it} - e^{-it}}{2i} \kildetag{Ligning 4.9}
\end{align*}

Lad $n \in \mathbb{N}$, da gælder
\begin{align*}
    \cos(n t) &= \Re((\cos(t)+\sin(t)i)^n) \\
    \sin(n t) &= \Im((\cos(t)+\sin(t)i)^n)
\end{align*}
\kilde{Sætning 4.5.1 - DeMoivres Formler}


\subsection{Polær form}
Et komplekst tal $z$ skrevet som \(|z| \cdot e^{i \cdot \arg(z)}\) er på polær form.
\[
    z = |z| \cdot \left(\cos(\arg(z)) + \sin(\arg(z)) \cdot i \right) = |z| \cdot e^{i \cdot \arg(z)} \kildetag{Definition 4.6.1}
\]

Generelt gælder at:
\begin{align*}
    |e^z| &= e^{\Re(z)} \\
    \arg(e^{z}) &= \Im(z)
\end{align*}
\kilde{Ligning 4.10}

Hvis $w \in \mathbb{C}$ og $w \neq 0$, så er alle $z$ der løse ligningen \(e^z = w\) givet ved:
\[
    z = \ln(|w|) + (\Arg(w) + p2\pi) i \quad p \in \mathbb{Z} \kildetag{Lemma 4.6.1}
\]
